\nuevaReceta{Arroz con Pulpo}{60 min}{receta:arroz_pulpo}

    \begin{minipage}[t]{0.35\textwidth}
        \section*{Ingredientes}
        \begin{itemize}[leftmargin=*]
            \item 350 g de arroz
            \item 500 g de pulpo cocido
            \item 1 Cebolla
            \item 1 Pimiento verde
            \item 1/2 Pimiento rojo
            \item 2 Dientes de ajo
            \item 2 Tomates maduros
            \item 1,2 litros de caldo de pescado
            \item 100 ml de vino blanco
            \item Pimentón dulce
            \item Azafrán o colorante
            \item Aceite de oliva
            \item Sal
        \end{itemize}
    \end{minipage}
    \hfill
    \begin{minipage}[t]{0.6\textwidth}
        \section*{Preparación}
        \begin{enumerate}
            \item \textbf{Preparar el pulpo:} Cortar el pulpo cocido en rodajas o trozos medianos. Reservar parte de las puntas de los tentáculos para decorar.
            \item \textbf{Hacer el sofrito:} Cocinar en una paellera con aceite la cebolla, los ajos y los pimientos bien picados.
            \item \textbf{Añadir el tomate:} Incorporar los tomates rallados y cocinar hasta que el sofrito pierda el agua y quede concentrado.
            \item \textbf{Incorporar el pulpo:} Añadir el pulpo y saltearlo durante unos minutos. Agregar el pimentón, remover rápidamente y verter el vino blanco.
            \item \textbf{Preparar el caldo:} Dejar que el vino reduzca y añadir el caldo caliente junto con el azafrán. Cocer durante 5 minutos y ajustar de sal.
            \item \textbf{Cocer el arroz:} Repartir el arroz por la paellera y cocinar unos 18 minutos, primero a fuego fuerte y después a fuego medio, sin remover.
            \item \textbf{Reposar:} Retirar del fuego cuando el arroz esté en su punto, colocar por encima el pulpo reservado y dejar reposar 5 minutos antes de servir.
        \end{enumerate}
    \end{minipage}

    \vspace{1em}
    \begin{tcolorbox}[colback=orange!5, colframe=orange!50, title=Consejo, size=small]
        Si has cocido el pulpo en casa, utiliza parte de su agua de cocción mezclada con el caldo de pescado para intensificar el sabor del arroz.
    \end{tcolorbox}
\end{receta}
